\chapter{Grid-Based PDE Solver}
\label{ch:pde}

The ODE approximation developed in Chapter~\ref{ch:ode} is computationally attractive: the Riccati system scales quadratically in the number of currencies and can be solved in milliseconds. However, this efficiency comes at a cost. The quadratic ansatz $\hat\theta(t,y) = -y^\T A(t)\,y - y^\T B(t) - C(t)$ restricts the value function to a particular parametric family, and the Hamiltonian is replaced by its second-order Taylor expansion around $p = 0$. At large inventories, where the true value function deviates from a quadratic and the momentum $p$ entering the Hamiltonian is far from zero, the approximation may break down. To assess when and how this happens, we turn to a direct numerical solution of the HJB equation.

In the two-currency case ($d = 2$), the inventory state is two-dimensional and the PDE is amenable to grid-based methods. We develop a fully implicit Euler solver with policy iteration (Howard's algorithm) for the HJB equation derived in Chapter~\ref{ch:model}. The key numerical challenge is the extreme stiffness introduced by the hedging Hamiltonian, which amplifies value function gradients by a factor of order $10^8$. We address this through a monotone finite difference scheme that guarantees the discrete comparison principle required for convergence. With this scheme in place, the solver handles the original model parameters without any artificial regularisation.

This chapter is organised as follows. Section~\ref{sec:pde-domain} describes the computational domain and grid. Section~\ref{sec:pde-spatial} presents the spatial discretisation of each term in the HJB equation. Section~\ref{sec:pde-stiffness} analyses the stiffness structure and explains why explicit time stepping is infeasible. Section~\ref{sec:pde-implicit} develops the fully implicit solver, including the policy iteration algorithm and the monotone differencing scheme. Section~\ref{sec:pde-convergence} discusses convergence behaviour and computational cost. Finally, Section~\ref{sec:pde-results} compares the PDE and ODE solutions at nominal parameters and across the parameter regimes identified in Chapter~\ref{ch:sensitivity}.


% ==================================================================
\section{Computational domain}
\label{sec:pde-domain}

We work with the two-currency model (USD, EUR) from Table~1 of \cite{barzykin2023}. Since USD is the reference currency with $\sigma_{\text{USD}} = 0$, the effective state variables are the two inventory components $y = (y_{\text{USD}}, y_{\text{EUR}})$, measured in millions of USD. The PDE is solved on a uniform Cartesian grid
\begin{equation}
  \label{eq:grid}
  \Omega = [-Y_{\max}, Y_{\max}]^2, \qquad
  y_i^{(k)} = -Y_{\max} + k\,\Delta y, \quad k = 0, 1, \ldots, n-1,
\end{equation}
with $Y_{\max} = 150$~M\$, $\Delta y = 1$~M\$, and $n = 301$ grid points per axis. This gives a total of $N = 301^2 = 90{,}601$ grid points.

Two considerations determine these choices. First, the grid spacing $\Delta y = 1$~M\$ is chosen so that every trade size in the model ($z \in \{1, 5, 10, 20, 50\}$~M\$) is an exact multiple of $\Delta y$. This ensures that the shifted point $y + z\,d_{ij}$ appearing in the quoting Hamiltonian always falls on a grid node, avoiding the need for interpolation. Second, the domain extends 50~M\$ beyond the region of interest $[-100, 100]$~M\$ in each direction. This buffer accommodates the largest trade size ($z = 50$~M\$): without it, the quoting Hamiltonian at $|y| > 50$~M\$ would lose the $z = 50$ contribution, artificially degrading the solution in the interior.

The terminal condition is $\theta(T, y) = -y^\T \kappa\, y$, which equals zero everywhere since the paper sets $\kappa = 0$. At the domain boundary $\partial\Omega$, we do not impose Dirichlet conditions. Instead, quoting contributions whose shifted lookup $y + z\,d_{ij}$ falls outside $\Omega$ are simply dropped, acting as an implicit inventory limit. For the implicit solver (Section~\ref{sec:pde-implicit}), boundary rows in the linear system are set to identity, effectively freezing boundary values. This simple treatment suffices because the 50~M\$ buffer ensures that the solution in the region of interest is insensitive to boundary effects.


% ==================================================================
\section{Spatial discretisation}
\label{sec:pde-spatial}

We write the HJB equation in backward time $\tau = T - t$, so that it takes the form
\begin{equation}
  \label{eq:hjb-backward}
  \partial_\tau \theta = \underbrace{\frac{1}{2}\operatorname{Tr}\!\bigl(\mathcal{D}(y)\,\Sigma\,\mathcal{D}(y)\,D^2_{yy}\theta\bigr)}_{\text{diffusion}}
  - \underbrace{\frac{\gamma}{2}\,y^\T \Sigma\, y}_{\text{running penalty}}
  + \underbrace{Q(y, \theta)}_{\text{quoting}}
  + \underbrace{\mathcal{H}_{\text{hedge}}(y, \nabla_y\theta)}_{\text{hedging}},
\end{equation}
where we have used that the drift $\mu = 0$ in the paper's example. The right-hand side defines the spatial operator $F[\theta]$, which we now discretise term by term.


% ------------------------------------------------------------------
\subsection{Diffusion}
\label{sec:pde-diffusion}

The diffusion term involves the Hessian of $\theta$ weighted by the state-dependent volatility. Expanding the trace,
\begin{equation}
  \label{eq:diffusion-expanded}
  \frac{1}{2}\operatorname{Tr}\!\bigl(\mathcal{D}(y)\,\Sigma\,\mathcal{D}(y)\,D^2_{yy}\theta\bigr)
  = \frac{1}{2}\sum_{i=1}^d \Sigma_{ii}\,y_i^2\,\frac{\partial^2\theta}{\partial y_i^2}
  + \sum_{i < j} \Sigma_{ij}\,y_i\,y_j\,\frac{\partial^2\theta}{\partial y_i\,\partial y_j}.
\end{equation}
Since the reference currency has $\sigma_{\text{USD}} = 0$, the entire first row and column of $\Sigma$ vanish. In the two-currency case, only a single diagonal term survives:
\begin{equation}
  \label{eq:diffusion-2ccy}
  \frac{1}{2}\,\sigma_{\text{EUR}}^2\,y_{\text{EUR}}^2\,\frac{\partial^2\theta}{\partial y_{\text{EUR}}^2}.
\end{equation}
The second derivative is approximated by the standard central difference stencil
\begin{equation}
  \label{eq:central-2nd}
  \frac{\partial^2\theta}{\partial y_i^2}\bigg|_k \approx \frac{\theta_{k+1} - 2\,\theta_k + \theta_{k-1}}{\Delta y^2},
\end{equation}
using one-sided (forward or backward) stencils at the domain boundary. The boundary stencils are first-order accurate rather than second-order, but they lie in the buffer zone where solution quality is not critical.


% ------------------------------------------------------------------
\subsection{Running penalty}
\label{sec:pde-penalty}

The running penalty $-(\gamma/2)\,y^\T \Sigma\, y$ does not involve $\theta$ and is precomputed once on the grid. In the two-currency case, the only nonzero contribution is
\begin{equation}
  -\frac{\gamma}{2}\,\sigma_{\text{EUR}}^2\,y_{\text{EUR}}^2.
\end{equation}
This term enters the right-hand side of~\eqref{eq:hjb-backward} as a source that penalises large inventory positions, driving the value function downward at the grid boundary.


% ------------------------------------------------------------------
\subsection{Quoting Hamiltonian}
\label{sec:pde-quoting}

The quoting term evaluates the exact Hamiltonian rather than the quadratic Taylor approximation used in the ODE. For each directed currency pair $(i,j)$, client tier $n$, and trade size $z_k$, the contribution at inventory $y$ is
\begin{equation}
  \label{eq:quoting-contribution}
  z_k\,\lambda_k\,H^n\!\left(\frac{\theta(y) - \theta(y + z_k\,d_{ij})}{z_k}\right),
\end{equation}
where $d_{ij} = e_i - e_j$ is the direction vector of the trade (the market maker receives currency~$i$ and delivers currency~$j$), $\lambda_k$ is the arrival rate for size $z_k$, and
\begin{equation}
  \label{eq:H-exact}
  H^n(p) = \sup_\delta\; f^n(\delta)\,(\delta - p), \qquad f^n(\delta) = \frac{1}{1 + e^{\alpha_n + \beta_n\,\delta}},
\end{equation}
is the quoting Hamiltonian for the logistic intensity function $f^n$. As shown in Chapter~\ref{ch:ode}, the supremum admits a semi-closed form via the Lambert~$W$ function:
\begin{equation}
  \label{eq:H-lambert}
  H^n(p) = \frac{w}{\beta_n}, \qquad w = W_0\!\bigl(e^{-(1 + \alpha_n + \beta_n p)}\bigr),
\end{equation}
where $W_0$ denotes the principal branch. The full quoting term $Q(y, \theta)$ is the sum of~\eqref{eq:quoting-contribution} over all directed pairs, tiers, and sizes. In the paper's model, the trade sizes are discrete ($z_k \in \{1, 5, 10, 20, 50\}$~M\$), so no numerical quadrature is required.

The argument of $H^n$ in~\eqref{eq:quoting-contribution} involves the finite difference $(\theta(y) - \theta(y + z_k\,d_{ij}))/z_k$, which couples each grid point to a shifted neighbor at distance $z_k$ along the direction $d_{ij}$. Since $\Delta y$ divides all trade sizes, these shifted points lie exactly on the grid. When the shifted point falls outside $\Omega$, the corresponding contribution is set to zero — the market maker does not quote that size at that inventory. For the two-currency case with two directed pairs, two tiers, and five trade sizes, this gives $2 \times 2 \times 5 = 20$ quoting contributions per grid point, each coupling a grid point to one shifted neighbor.


% ------------------------------------------------------------------
\subsection{Hedging Hamiltonian}
\label{sec:pde-hedging}

For each canonical currency pair $(i,j)$ with $i < j$, the hedging Hamiltonian evaluates
\begin{equation}
  \label{eq:H-hedge}
  \mathcal{H}^{i,j}(p) = \sup_\xi\;\bigl(p\,\xi - \psi^{i,j}\,|\xi| - \eta^{i,j}\,\xi^2\bigr)
  = \frac{\bigl(\max(|p| - \psi^{i,j},\; 0)\bigr)^2}{4\,\eta^{i,j}},
\end{equation}
where $\psi^{i,j}$ and $\eta^{i,j}$ are the proportional spread cost and quadratic market impact for the D2D pair, and the argument is
\begin{equation}
  \label{eq:p-hedge}
  p^{i,j}(y) = \frac{\partial\theta}{\partial y_i} - \frac{\partial\theta}{\partial y_j}
  + k_i\,y_i\!\left(1 + \frac{\partial\theta}{\partial y_i}\right) - k_j\,y_j\!\left(1 + \frac{\partial\theta}{\partial y_j}\right).
\end{equation}
In the paper's example, the market impact parameters $k_i$ are negligibly small, so $p$ reduces to the difference of value function gradients along the two currency axes. The gradients are approximated by central finite differences in the grid interior.

The optimal hedge rate at the supremum is
\begin{equation}
  \label{eq:xi-star}
  \xi^*(p) = \begin{cases}
    (p - \psi) / (2\eta) & \text{if } p > \psi, \\
    (p + \psi) / (2\eta) & \text{if } p < -\psi, \\
    0 & \text{if } |p| \leq \psi.
  \end{cases}
\end{equation}
The dead zone $|p| \leq \psi$ reflects the fact that the proportional spread cost makes small hedges unprofitable. Outside the dead zone, the optimal hedge rate is proportional to $1/(2\eta)$, which is of order $5 \times 10^8$ for the paper's parameters. This extreme sensitivity to the value function gradient is the dominant source of numerical stiffness, as we discuss in the next section.

Once the optimal hedge rate $\xi^*$ is computed, the hedging contribution to the right-hand side of~\eqref{eq:hjb-backward} can be decomposed into a term involving $\xi^*\,\nabla_y\theta$ (an advection-like coupling to the value function) and a source term collecting the execution costs $-\psi\,|\xi^*| - \eta\,(\xi^*)^2$. This decomposition is central to the linearised system in the implicit solver.


% ==================================================================
\section{Stiffness and the need for implicit time stepping}
\label{sec:pde-stiffness}

The simplest approach to marching~\eqref{eq:hjb-backward} forward in $\tau$ is explicit Euler:
\begin{equation}
  \label{eq:explicit-euler}
  \theta^{m+1} = \theta^m + \Delta t\; F[\theta^m],
\end{equation}
where $F$ is the spatial operator. Explicit time stepping has the virtue of simplicity — each step requires only one evaluation of $F$ — but it is subject to stability constraints.

The diffusion term imposes the standard CFL condition $\Delta t < \Delta y^2 / (\sigma^2\,Y_{\max}^2)$. With the paper's parameters ($\sigma_{\text{EUR}} = 80$~bps, $Y_{\max} = 150$~M\$, $\Delta y = 1$~M\$), this yields $\Delta t < 0.69$~days, which is very generous relative to the time horizon $T = 0.05$~days. The diffusion is not the binding stability constraint.

The hedging Hamiltonian, however, introduces severe stiffness. The gain $1/(4\eta) \approx 2.5 \times 10^8$ in~\eqref{eq:H-hedge} means that even modest value function gradients produce enormous Hamiltonian values, creating a positive feedback loop: a small perturbation in $\theta$ amplifies the gradient, which amplifies $\mathcal{H}_{\text{hedge}}$, which amplifies $\theta$ further. This mechanism causes numerical blow-up regardless of how small $\Delta t$ is chosen, because the nonlinear Hamiltonian does not satisfy a standard linear CFL condition.

An intermediate approach --- treating the diffusion term implicitly while keeping the Hamiltonians explicit (an IMEX or semi-implicit scheme) --- was also considered. Since the stability constraint is dominated by the hedging Hamiltonian rather than the diffusion, this scheme requires the same artificial scaling of $\eta$ as the explicit method and offers no improvement. The stiffness can only be resolved by treating the hedging term implicitly, which leads naturally to a fully implicit scheme.


% ==================================================================
\section{Fully implicit solver with policy iteration}
\label{sec:pde-implicit}


% ------------------------------------------------------------------
\subsection{Implicit Euler formulation}
\label{sec:pde-implicit-euler}

The fully implicit Euler scheme reads
\begin{equation}
  \label{eq:implicit-euler}
  \theta^{m+1} - \Delta t\; F[\theta^{m+1}] = \theta^m.
\end{equation}
This is a nonlinear equation for $\theta^{m+1}$, because the spatial operator $F$ contains the quoting and hedging Hamiltonians, which involve suprema over controls. In return for this added complexity, the implicit scheme avoids the CFL restriction: the stiff hedging term is absorbed into the left-hand side and handled by the linear algebra solver.


% ------------------------------------------------------------------
\subsection{Policy iteration}
\label{sec:pde-pi}

We solve the nonlinear equation~\eqref{eq:implicit-euler} at each time step using policy iteration (Howard's algorithm), which alternates between two steps:

\begin{enumerate}
  \item \textbf{Fix the controls.} Given a current estimate $\theta^k$ (initialised to $\theta^m$), compute the optimal quotes $\delta^*(y)$ and hedge rates $\xi^*(y)$ at every grid point using the formulas~\eqref{eq:H-lambert} and~\eqref{eq:xi-star}.

  \item \textbf{Solve the linearised system.} With the controls frozen, the suprema in $F$ are removed and the operator becomes linear in $\theta^{k+1}$. This yields a sparse linear system
  \begin{equation}
    \label{eq:linearised-system}
    (I - \Delta t\,A)\;\theta^{k+1} = \theta^m + \Delta t\; s,
  \end{equation}
  where $A$ is a sparse matrix encoding how each grid point's value depends on its neighbours through finite differences, and $s$ is a source vector collecting all terms that do not depend on $\theta^{k+1}$.
\end{enumerate}

The iteration terminates when the relative change falls below a prescribed tolerance:
\begin{equation}
  \frac{\|\theta^{k+1} - \theta^k\|_\infty}{\max(1,\;\|\theta^{k+1}\|_\infty)} < \varepsilon_{\text{PI}}.
\end{equation}
Policy iteration is equivalent to Newton's method applied to the Bellman equation and inherits its quadratic convergence rate, typically requiring only a small number of iterations per time step.


% ------------------------------------------------------------------
\subsection{Assembly of the linearised system}
\label{sec:pde-assembly}

With the controls frozen, each term in the spatial operator~\eqref{eq:hjb-backward} contributes either to the sparse matrix $A$ (through its dependence on $\theta^{k+1}$) or to the source vector $s$ (through quantities known from the current iterate $\theta^k$). We now describe these contributions.

The \emph{running penalty} $-(\gamma/2)\,y^\T \Sigma\,y$ has no $\theta$ dependence and enters entirely as a source term. It is precomputed once on the grid.

The \emph{diffusion} term is linear in $\theta$ and enters entirely into $A$. For each axis $i$ with $\Sigma_{ii} > 0$, the central difference stencil~\eqref{eq:central-2nd} contributes three entries per row: $\theta_{k-1}$, $\theta_k$, and $\theta_{k+1}$ along axis $i$, weighted by $\frac{1}{2}\,\Sigma_{ii}\,y_i^2 / \Delta y^2$.

The \emph{quoting Hamiltonian} requires more care. With the optimal markup $\delta^*$ and fill probability $f^*$ frozen, the contribution of a single quoting event $(n, i, j, z_k)$ at grid point $y$ is
\begin{equation}
  z_k\,\lambda_k\,f^*\!\left(\delta^* - \frac{\theta(y) - \theta(y + z_k\,d_{ij})}{z_k}\right)
  = \underbrace{z_k\,\lambda_k\,f^*\,\delta^*}_{\to\; s}
  \;-\; \underbrace{\lambda_k\,f^*\,\bigl(\theta(y) - \theta(y + z_k\,d_{ij})\bigr)}_{\to\; A}.
\end{equation}
The first part is a source term (the expected profit per quoting event). The second part contributes two entries per row of $A$: a diagonal entry $-\lambda_k\,f^*$ and an off-diagonal entry $+\lambda_k\,f^*$ at the shifted position $y + z_k\,d_{ij}$. With 20 quoting contributions in the two-currency case, this adds up to 40 off-diagonal entries per row.

The \emph{hedging Hamiltonian} with frozen hedge rate $\xi^*$ decomposes into a source term $-\psi\,|\xi^*| - \eta\,(\xi^*)^2$ (the execution cost) and a linear term $\xi^*\,(\partial\theta/\partial y_i - \partial\theta/\partial y_j)$ (the advection coupling). The advection term contributes to $A$ through the gradient stencil, adding up to four entries per row (two per axis). The discretisation of this gradient is the subject of the next subsection.

In total, the matrix $I - \Delta t\,A$ has approximately 48 nonzero entries per row, giving roughly $4.3 \times 10^6$ nonzeros for the $90{,}601 \times 90{,}601$ system. The matrix is assembled in COO format and converted to CSR for the sparse direct solve.


% ------------------------------------------------------------------
\subsection{Monotone differencing and the M-matrix property}
\label{sec:pde-monotone}

The hedging advection term $\xi^*\,\partial\theta/\partial y_i$ is a first-order transport-like operator with velocity $v = \xi^*(y)$. For the paper's parameters, $|\xi^*|$ can reach $50{,}000$~M\$/day outside the dead zone, producing matrix entries of magnitude $\Delta t\,|\xi^*|/\Delta y \approx 60$ when central differences are used for the gradient. Since the diagonal entry from this term is zero under central differencing, and the identity contributes only~1 to the diagonal, the matrix $I - \Delta t\,A$ loses diagonal dominance. The resulting linear system is ill-conditioned, and the sparse direct solver produces erroneous solutions that cause policy iteration to diverge.

The remedy is to use directional (upwind) differences for the hedging gradient, chosen so that the matrix $I - \Delta t\,A$ is an \emph{M-matrix}: a matrix with non-positive off-diagonal entries and a dominant positive diagonal. Specifically, for the advection operator $+v\,\partial\theta/\partial y_i$ entering the right-hand side of~\eqref{eq:hjb-backward}, we discretise as
\begin{equation}
  \label{eq:upwind}
  v\,\frac{\partial\theta}{\partial y_i}\bigg|_k \approx \begin{cases}
    v\,(\theta_{k+1} - \theta_k) / \Delta y & \text{if } v > 0, \\
    v\,(\theta_k - \theta_{k-1}) / \Delta y & \text{if } v < 0.
  \end{cases}
\end{equation}
Under this choice, the contributions to $I - \Delta t\,A$ from the hedging gradient are:
\begin{alignat}{2}
  &\text{diagonal:}      &\quad &{+}\;\Delta t\,|v|/\Delta y, \label{eq:upwind-diag} \\
  &\text{off-diagonal:}  &\quad &{-}\;\Delta t\,\max(v, 0)/\Delta y \;\;\text{(at $k+1$)}, \notag \\
  &                      &\quad &{+}\;\Delta t\,\min(v, 0)/\Delta y \;\;\text{(at $k-1$)}. \notag
\end{alignat}
The diagonal entry is always positive, and both off-diagonal entries are non-positive, regardless of the sign or magnitude of $v$. Combined with the positive contributions from the identity and the diffusion operator, this guarantees that $I - \Delta t\,A$ is an M-matrix.

We note that the upwind direction in~\eqref{eq:upwind} is the opposite of the standard convention for the transport equation $\partial_\tau u + v\,\partial u/\partial y = 0$, where $v > 0$ calls for a backward difference. The difference arises because the advection term in~\eqref{eq:hjb-backward} enters with a positive sign on the right-hand side: positive $v\,\partial\theta/\partial y$ makes $\theta$ grow, and the implicit scheme must move this growth to the left-hand side with the correct sign to maintain the M-matrix structure.

The M-matrix property is significant for two reasons. First, it ensures a discrete comparison principle: if two solutions satisfy an ordering at the boundary and in the source, the ordering is preserved in the interior. This is the discrete analogue of the maximum principle for elliptic PDEs, and it is a necessary condition for the convergence of monotone numerical schemes in the framework of Barles and Souganidis~\cite{barles1991}. Second, it guarantees the convergence of policy iteration, which relies on the monotonicity of the discrete Bellman operator. Without the M-matrix structure, policy iteration may oscillate or diverge even when the continuous problem is well-posed.

The price of monotone differencing is a reduction in spatial accuracy from second order (central differences) to first order. In practice, this loss is confined to the hedging gradient and is most noticeable at the grid boundary, where the first-order stencil interacts with the boundary treatment. In the interior of the grid, where the solution is smooth and the region of interest lies, the first-order error is small relative to the time discretisation error. The central difference stencils for the diffusion operator and the quoting Hamiltonian are unaffected.


% ==================================================================
\section{Convergence and computational cost}
\label{sec:pde-convergence}

With the monotone differencing scheme in place, the solver handles the paper's original parameters ($\eta = 10^{-9}$ in decimal units) without any artificial scaling or continuation. A tolerance of $\varepsilon_{\text{PI}} = 10^{-4}$ is used for the policy iteration. Detailed iteration traces reveal the following convergence pattern:

At the first time step (cold start from $\theta(T) = 0$), the policy iteration requires approximately 7 inner iterations. The relative change drops rapidly from $3.6 \times 10^{-2}$ to $2.8 \times 10^{-4}$ over the first five iterations, with the value at the origin $\theta(0,0)$ stabilising after iteration~5. The residual then stalls at approximately $2.5 \times 10^{-5}$, oscillating without further decrease. This stall is caused by the first-order accuracy of the upwind stencil at boundary points and does not affect the solution in the region of interest. Setting $\varepsilon_{\text{PI}} = 10^{-4}$, safely above the stall level, allows the iteration to terminate promptly.

At subsequent time steps, the previous solution provides an effective warm start, and the policy iteration converges in 2--4 inner iterations. With 20 time steps over the horizon $T = 0.05$~days, the solver completes in approximately 9~minutes on a $301 \times 301$ grid, with each sparse linear solve taking roughly 0.1~seconds via a direct sparse factorisation.

We remark that a continuation method in $\eta$ (solving a sequence of problems with decreasing $\eta$, using each solution to warm-start the next) was also tested. The monotone scheme alone was found to be sufficient for convergence at the original $\eta$, and the continuation added computational cost without improving the solution. This is consistent with the analysis in Section~\ref{sec:pde-monotone}: the stability issue is not the magnitude of $\eta$ per se, but the loss of the M-matrix property under central differencing.


% ==================================================================
\section{Numerical results}
\label{sec:pde-results}

We now compare the PDE solution against the ODE approximation from Chapter~\ref{ch:ode}. The comparison focuses on the three quantities of interest introduced in Chapter~\ref{ch:sensitivity}: the neutral spread $\delta^*(y = 0)$, the inventory skew $\delta^*(y = 10) - \delta^*(y = 0)$, and the hedge rate $\xi^*(y = 10)$.


% ------------------------------------------------------------------
\subsection{Nominal parameters}
\label{sec:pde-nominal}

At the paper's nominal parameter values, the PDE and ODE value functions agree closely. Along inventory slices through the origin, the two solutions are visually indistinguishable over the range $|y| \leq 100$~M\$. The maximum absolute difference is approximately $1.1 \times 10^{-4}$ in the value function, occurring at $|y_{\text{EUR}}| \approx 75$~M\$, compared to a typical value of $\theta(0,0) \approx 1.06 \times 10^{-2}$. This corresponds to a relative discrepancy of roughly 1\%.

The PDE value function is systematically higher than the ODE approximation at large inventories. The ODE relies on two coupled simplifications: a quadratic Taylor expansion $\hat{H}(p) \approx \alpha_0 + \alpha_1\,p + \frac{1}{2}\,\alpha_2\,p^2$ of the Hamiltonian and a quadratic ansatz $\hat\theta = -y^\T A\,y - y^\T B - C$ for the value function. Their combined effect is to slightly underestimate the true value at large $|y|$ --- the PDE, which captures the full nonlinear structure of both $H$ and $\theta$, shows that the market maker extracts marginally more value from quoting at extreme inventories than the ODE predicts.

% TODO: Add figure — PDE vs ODE value function along EUR inventory slice
% \begin{figure}[ht]
%   \centering
%   \includegraphics[width=0.8\textwidth]{figures/pde_ode_value_function.pdf}
%   \caption{Value function $\theta(0, y)$ along the EUR inventory axis ($y_{\text{USD}} = 0$) at $t = 0$, computed by the PDE solver (solid) and the ODE approximation (dashed). The inset shows the difference $\theta_{\text{PDE}} - \theta_{\text{ODE}}$.}
%   \label{fig:pde-ode-value}
% \end{figure}

The three quantities of interest confirm this picture. At the origin, the neutral spread and inventory skew from the PDE match those from the ODE to within the resolution of the grid. The hedge rate, which depends on the value function gradient rather than its level, also agrees well. These results confirm that the ODE approximation is accurate for the paper's nominal parameter set, consistent with the brief validation reported in~\cite{barzykin2023}.


% ------------------------------------------------------------------
\subsection{Sensitivity across parameter regimes}
\label{sec:pde-sensitivity}

The sensitivity analysis in Chapter~\ref{ch:sensitivity} identified four parameters that drive the quantities of interest: $\sigma_{\text{EUR}}$, $\gamma$, $\lambda_{\text{scale}}$, and $\eta$. To assess whether the ODE approximation remains reliable across their uncertainty ranges, we solve the PDE at the low and high end of each parameter's range (Table~\ref{tab:sensitivity-ranges} in Chapter~\ref{ch:sensitivity}), with all other parameters held at their nominal values. This produces 9~PDE solves in total: one at the nominal point and two for each of the four parameters.

% TODO: Add table — PDE vs ODE QoIs at each parameter configuration
% \begin{table}[ht]
%   \centering
%   \caption{Quantities of interest from the PDE and ODE solvers at nominal parameters and at the edges of the sensitivity ranges for $\sigma$, $\gamma$, $\lambda$, and $\eta$.}
%   \label{tab:pde-ode-qois}
%   \begin{tabular}{l rrr rrr}
%     \toprule
%     & \multicolumn{3}{c}{ODE} & \multicolumn{3}{c}{PDE} \\
%     \cmidrule(lr){2-4} \cmidrule(lr){5-7}
%     Config & $\delta^*(0)$ & skew & $\xi^*$ & $\delta^*(0)$ & skew & $\xi^*$ \\
%     \midrule
%     nominal & & & & & & \\
%     $\sigma$ low & & & & & & \\
%     $\sigma$ high & & & & & & \\
%     \ldots & & & & & & \\
%     \bottomrule
%   \end{tabular}
% \end{table}

For each parameter configuration, we compare the PDE and ODE values of the three quantities of interest, as well as the sensitivity magnitudes (the change in each QoI between the low and high parameter values). This comparison tests two things: whether the ODE accurately captures the absolute level of each QoI across parameter regimes, and whether it correctly predicts the direction and magnitude of the sensitivity.

% TODO: Add figure — bar chart comparing ODE vs PDE sensitivities
% \begin{figure}[ht]
%   \centering
%   \includegraphics[width=\textwidth]{figures/pde_ode_sensitivity_comparison.pdf}
%   \caption{Sensitivity of the three quantities of interest to the four key parameters, as computed by the PDE solver and the ODE approximation. Each bar shows the difference in the QoI between the high and low parameter values.}
%   \label{fig:pde-ode-sensitivity}
% \end{figure}


% ------------------------------------------------------------------
\subsection{Discussion}
\label{sec:pde-discussion}

The comparison across parameter regimes serves two purposes within the thesis. First, it validates the PDE solver itself: if the PDE and ODE solutions agree in regimes where the quadratic approximation is expected to be accurate, this provides confidence that the spatial operators and the implicit scheme are implemented correctly. Second, it validates the ODE-based sensitivity conclusions from Chapter~\ref{ch:sensitivity}: if the PDE confirms the same sensitivity rankings and magnitudes, it justifies using the computationally cheap ODE for the full uncertainty quantification study, including potential extensions to more than two currencies where the PDE becomes infeasible.

Any systematic discrepancies between the PDE and ODE solutions are most likely to appear in two regimes. First, at high volatility or high risk aversion, the optimal inventory skew increases, pushing the value function further from the quadratic form and the momentum $p$ further from zero, where the Taylor expansion of $H$ is least accurate. Second, at low $\eta$, the hedging Hamiltonian's gain $1/(4\eta)$ is largest, amplifying any differences in the value function gradient. In both cases, the PDE provides the more reliable benchmark, and the magnitude of the discrepancy characterises the domain of validity of the ODE approximation.
